\documentclass[12pt,letterpaper]{article}
\usepackage[utf8]{inputenc}
\usepackage[spanish]{babel}
\usepackage{times}
\usepackage[margin=2.54cm]{geometry}
\usepackage{setspace}
\usepackage{titlesec}
\usepackage{enumitem}
\usepackage{hyperref}
\usepackage{fancyhdr}
\usepackage{microtype}
\sloppy

\setlength{\parindent}{0pt}
\setlength{\parskip}{6pt}
\singlespacing

\pagestyle{fancy}
\fancyhf{}
\fancyfoot[C]{\thepage}
\renewcommand{\headrulewidth}{0pt}

\titleformat{\section}{\bfseries}{}{0em}{}
\titlespacing{\section}{0pt}{6pt}{2pt}

\hypersetup{colorlinks=true, urlcolor=blue, linkcolor=black}

\begin{document}

\begin{center}
{\bfseries DETRÁS DEL MECANISMO DE CONSULTA }\\[4pt]
\textit{K.V. Ramirez Ambrocio}\\
7690-22-2239 \quad Universidad Mariano Gálvez\\
Seminario de Tecnologías de Información\\
kramireza10@miumg.edu.gt
\end{center}

\vspace{6pt}

\section*{Resumen}
Este ensayo busca explicar qué hace que una encuesta funcione bien como herramienta para recolectar datos, sobre todo cuando se usa en proyectos de investigación tecnológica o de seminario. Para esto se revisó artículos sobre metodología de encuestas, principalmente de la Asociación Americana de Investigación de Opinión Pública (AAPOR) y de estudios publicados en Public Opinion Quarterly, comparando esa información con la guía metodológica del curso. Como resultado se encontraron cinco puntos clave para que una encuesta funcione mejor: tener bien claro el objetivo, redactar bien las preguntas, elegir correctamente a quién se le va a preguntar, hacer una prueba piloto antes de aplicarla y controlar el error total de encuesta, que es la suma de los errores de cobertura, muestreo, no respuesta y medición. La conclusión principal es que una encuesta buena no es la que junta más respuestas, sino la que consigue datos de calidad y representativos, por lo que hay que planearla con el mismo cuidado que cualquier otro instrumento de investigación.

\textbf{Palabras clave:} encuesta, diseño de instrumentos, error total de encuesta, muestreo, recolección de datos.

\section*{Desarrollo del tema}

Una encuesta, básicamente, es una forma de conseguir información directamente de un grupo de personas para saber qué opinan, cómo se comportan o qué características tienen. Aunque parezca algo fácil de armar, los artículos sobre el tema coincide en que los resultados dependen de decisiones que se toman mucho antes de escribir la primera pregunta. Antes de ponerse a redactar preguntas, hay que tener bien claro qué se quiere averiguar y qué problema se está tratando de resolver, porque este paso muchas veces se salta y termina generando encuestas desordenadas y poco útiles.

Ese primer paso, tener claro el objetivo, es el que determina todo lo demás. Si no se sabe con claridad qué se quiere medir, es imposible saber qué preguntas hacer, a quién dirigir la encuesta. Por eso, antes de armar la encuesta conviene escribir en una sola oración cuál es la meta de la encuesta, y de ahí sacar cada pregunta, quitando las que no aporten directamente a esa meta.

El segundo punto importante es cómo se redactan las preguntas. Aquí los errores más comunes son las preguntas dobles (que en realidad están preguntando dos cosas al mismo tiempo), las preguntas que ya sugieren una respuesta, y el uso de palabras poco claras. Por ejemplo, si se pregunta "¿qué tan satisfecho está con la calidad y el precio del producto?", se están mezclando dos cosas distintas y la persona no va a saber a cuál responder; lo correcto es separar esa pregunta en dos. También se recomienda organizar las preguntas usando un filtro, empezar con preguntas generales y terminar con las más específicas, para no aburrir a la persona desde el inicio.

El tercer punto, es cómo se elige la muestra. No basta con mandarle la encuesta a mucha gente, lo que importa de verdad es que esa muestra represente bien a la población que se quiere estudiar. Aquí entra el concepto de \textit{error total de encuesta}, que es básicamente un marco que junta todas las fuentes de error que pueden afectar qué tan precisos son los resultados. Este marco diferncia 4 tipos de error: el error de cobertura, que pasa cuando la muestra no incluye a toda la población que interesa; el error de muestreo, que es la variación normal que existe por trabajar con una muestra y no con toda la población; el error de no respuesta, que aparece cuando parte de la gente seleccionada no contesta, y esas personas pueden ser distintas a las que sí respondieron; y el error de medición, que tiene que ver con cómo se formulan y se aplican las preguntas.

Entender estos 4 errores sirve porque deja claro que una encuesta efectiva no es la que junta más respuestas, sino la que logra reducir estos errores dentro del tiempo y presupuesto disponibles. Por ejemplo, de nada sirve invertir en tener más gente encuestada si el la muestra está dejando fuera a un grupo importante de la población, porque ahí el problema es de cobertura, no de tamaño de muestra.

El cuarto punto es la prueba piloto. Antes de aplicar la encuesta a todo el mundo, se recomienda probarla primero con un grupo pequeño de personas para detectar preguntas confusas, errores de redacción o problemas técnicos. Este paso, aunque casi siempre se salta por falta de tiempo, es el que permite corregir errores de diseño antes de que afecten a toda la muestra, y se considera una práctica estándar en investigación de opinión pública.

Por último está el control de calidad durante la recolección de los datos. Esto incluye revisar que la muestra se esté seleccionando bien, que el todo esté funcionando sin errores técnicos y que los datos se estén procesando de forma correcta. En el caso de encuestas digitales, herramientas como cambiar el orden de las preguntas al azar o usar lógica de salto ayudan a reducir problemas y a mantener el interés de la persona hasta terminar la encuesta .

Aquí es importante saber diferenciar entre información que está bien respaldada y contenido que es más como falsos positivos. Mucho de lo que se encuentra en internet sobre cómo hacer la encuesta perfecta viene de plataformas publicitadas de encuestas, que lo que buscan es promover su propia aplicación, esas fuentes dan consejos prácticos que sirven, pero no siempre tienen evidencia detrás. En cambio, el error total de encuesta vienen de artpiculos académicos revisada por otros investigadores y dan mayor respaldo para justificar decisiones de diseño en un trabajo de investigación.

\section*{Observaciones y comentarios}

Al revisar los articulos sobre este tema, algo que llamó la atención es que la mayoría de las guías  que hay en internet se enfocan casi solo en cómo redactar las preguntas, dejando de lado cosas igual de importantes como el diseño de la muestra y el control del error de no respuesta. Esto puede hacer que alguien diseñe una encuesta con preguntas muy bien escritas, pero aplicada a una muestra poco representativa, lo cual daña los resultados sin que la persona se dé cuenta del problema real.

Otra cosa que se nota es que el error total de encuesta, aunque es el concepto más riguroso desde el punto de vista académico, es también el menos conocido fuera de los círculos de investigación de opinión pública. Sería bueno que en los cursos de metodología de investigación este marco se enseñara junto a como redactar preguntas, ya que uno complementa al otro y no lo reemplaza.

Como problema de este analisis, hay que mencionar que no se comparó entre distintas formas de aplicar la encuesta (en línea, telefónica, presencial), porque cada una tiene sus propias fuentes de error y ese tema solo ya daría para un ensayo completo. Como posible investigación futura, sería interesante estudiar cómo el orden al azar y la lógica de salto en encuestas digitales afectan al error de medición, en comparación con las encuestas en papel.

\section*{Conclusiones}

\begin{enumerate}[label=\arabic*., leftmargin=*]
\item Una encuesta efectiva se realiza desde que se define bien el objetivo de la investigación, no desde que se redactan las preguntas  sin un objetivo claro, ninguna técnica de diseño que se use después corrige el problema.
\item Lo que hace buena a una encuesta depende más de qué tan representativa es la muestra y de controlar los errores de cobertura y no respuesta, que de cuánta gente en total contestó.
\item El error total de encuesta da una base académica importante para justificar decisiones metodológicas, y es más confiable que la mayoría de guías comerciales que hay en internet.
\item La prueba piloto es un paso que casi siempre se omite, pero es la forma más directa de detectar errores de medición antes de aplicar el instrumento a toda la muestra.
\item Hay que diferenciar entre las recomendaciones prácticas de plataformas comerciales de encuestas y los hallazgos respaldados por articulos académica revisada por pares, porque ambas fuentes cumplen funciones distintas dentro del proceso de diseño.
\end{enumerate}

\section*{Referencias}

\begin{spacing}{1.0}
\setlength{\parindent}{-1.27cm}
\setlength{\leftskip}{1.27cm}

Asociación Americana de Investigación de Opinión Pública. (2024). \textit{Best practices for survey research}. AAPOR. \url{https://aapor.org/standards-and-ethics/best-practices/}

Biemer, P. P. (2010). Total survey error: Design, implementation, and evaluation. \textit{Public Opinion Quarterly, 74}(5), 817--848. \url{https://doi.org/10.1093/poq/nfq058}

Dynata. (2025, 3 de octubre). \textit{Survey design best practices: Designing reliable surveys}. \url{https://www.dynata.com/why-dynata/resources/blog/survey-design-best-practices/}

Groves, R. M., \& Lyberg, L. (2010). Total survey error: Past, present, and future. \textit{Public Opinion Quarterly, 74}(5), 849--879. \url{https://doi.org/10.1093/poq/nfq065}

\end{spacing}

\vspace{12pt}
\noindent\small Repositorio Git: \url{https://github.com/Valki11/seminario-encuesta.git}

\noindent\small Encuesta: \url{https://forms.gle/9YTJLrA2NbPkcdGCA}

\end{document}